\documentclass[aspectratio=169]{beamer}
\usepackage[T1]{fontenc}
\usepackage[utf8]{inputenc}
\usepackage[ngerman,english]{babel}
\usepackage{texstage}
\usepackage{texstage-fd-lab}
\setbeamertemplate{navigation symbols}{}
\title{What follows from a dependency?}
\author{TeXStage Functional Dependencies Lab}
\date{}
\begin{document}
\begin{frame}\titlepage
\small Use given business rules to derive attributes, identify minimal keys, and prove or refute dependencies.
\end{frame}
\begin{frame}{What does a workshop determine?}
\fdlab{closure}{.96\linewidth}{.78\textheight}
\end{frame}
\begin{frame}{Start with a person and a workshop}
\fdlab{practice}{.96\linewidth}{.78\textheight}
\end{frame}
\begin{frame}{Identify one workshop registration}
\fdlab{keys_one}{.96\linewidth}{.78\textheight}
\end{frame}
\begin{frame}{Find every candidate key}
\fdlab{keys_all}{.96\linewidth}{.78\textheight}
\end{frame}
\begin{frame}{Does a workshop determine its building?}
\fdlab{implied}{.96\linewidth}{.78\textheight}
\end{frame}
\begin{frame}{Does a tutor determine the workshop?}
\fdlab{refuted}{.96\linewidth}{.78\textheight}
\end{frame}
\begin{frame}{Minimal does not mean smallest}
\fdlab{different_sizes}{.96\linewidth}{.78\textheight}
\end{frame}
\begin{frame}{Can the empty set be a key?}
\fdlab{empty_key}{.96\linewidth}{.78\textheight}
\end{frame}
\begin{frame}{What have we shown?}
\begin{itemize}
\item An FD applies only when all attributes on its left are known.
\item Repeating productive steps reaches a fixed point: the complete closure.
\item A candidate key is minimal by inclusion; keys may have different sizes.
\item An implication follows by derivation, or fails by a valid counterexample.
\item Uniqueness in one sample table does not establish a schema constraint.
\end{itemize}
\end{frame}
\end{document}
